\chapter{バリオン生成、宇宙の滑らかさ、ハッブル摩擦}

\lectureoverview{二つの小さな無次元数が今夜の講義を整理する。第一は生き残った物質--反物質の超過であり、講義では\(10^8\)分の1に丸められる。その力学的説明にはバリオン数の破れ、粒子--反粒子の非対称性、熱平衡からのずれ、すなわちサハロフの三条件が必要である。第二は原始密度揺らぎであり、脱結合時で\(10^{-5}\)の程度である。重力はそのようなさざ波を消すのではなく増幅するので、その小ささがインフレーションを動機づける。次に我々は、膨張宇宙における一様スカラー場の運動方程式を導き、\(3H\dot\phi\)をハッブル摩擦として同定することにより、インフレーションの力学を準備する。}

\section{逆さまにされたエントロピー問題}

今夜インフレーションを始めるが、その前にバリオン生成、すなわち反物質に対する物質の超過の生成を終えなければならない。この枠組みは、いかなる詳細な数値模型よりも確固としている。だからこそ、その中心的な数が第一原理から導かれていないにもかかわらず、この主題は簡潔に要約できるのである。

この問題への歴史的な道は1980年に始まり、ボブ・ワゴナーがサバス・ディモプロスとレナード・サスキンドに、なぜ宇宙のエントロピーはこれほど大きいのかと尋ねた。宇宙論的に言えば、この問いは、生き残った陽子一つあたりなぜこれほど多くの光子が存在するのかということであった。

非常に高温では、光子、クォーク、反クォーク、電子、陽電子がすべて豊富に存在し、熱平衡に近かった。クォークと反クォークはそれぞれ、光子密度と同程度の広い桁の数密度を持っていた。この観察はすでに、問いを逆にすることを示唆する。なぜ対消滅は、まったく対応するもののない物質を残したのか。

この構造的議論の優先権はアンドレイ・サハロフに帰属する。ディモプロスとサスキンドは1980年に独立に密接に関連する条件に到達し、その後、歴史的寄稿を準備する中でサハロフの1967年の論文を見つけた。これらの条件は今や正当にもサハロフの名で知られている。

エントロピーとの関連をここで明示できる。黒体放射では、エントロピーは1の程度の固定因子を除いて光子数に比例する。実際、
\begin{equation}
s_\gamma=\frac{4\rho_\gamma}{3T},
\qquad
n_\gamma\propto T^3,
\end{equation}
したがって共動する熱的光子気体は\(S_\gamma\propto N_\gamma\)を持つ。数値定数が問題なのではない。謎は、光子と陽子の巨大な比である。

\begin{classroomqa}
\audiencequestion{光子の存在量はエントロピーとどのように関係するのか。}
\lecturerresponse{黒体放射は熱的であり、そのエントロピーは平均光子数に比例する。この文脈では、なぜ宇宙のエントロピーが大きいのかと問うことは、本質的に、なぜ宇宙マイクロ波背景放射が陽子一つあたりこれほど多くの光子を含むのかと問うことである。比例関係には数値定数が含まれるが、それらは桁の謎には無関係である。\lecturetimestamp{00:04:49--00:05:55}}
\end{classroomqa}

宇宙が冷えるにつれて粒子と反粒子は対消滅し、放射は残って最終的に脱結合した。したがって驚くべき事実は、光子が生き残ったことではなく、対消滅がごくわずかな対応するもののない物質の集団を残したことである。

したがって有用な無次元数は
\begin{equation}
\eta_B
\equiv
\frac{N_B-N_{\bar B}}{N_\gamma}
\sim
\frac{N_p-N_{\bar p}}{N_p+N_{\bar p}}
\sim 10^{-8},
\label{eq:c08-baryon-asymmetry}
\end{equation}
であり、最後の値は講義で用いられる丸めた桁であり、精密な決定値ではない。一見奇妙な数が現れると、``単にそう始まった''という説明は満足のいく終着点ではなくなる。定量的な機構は、なぜ答えが小さいのか、そしてなぜそれが確定した符号を持つのかを説明すべきである。

余剰を\emph{物質}（\emph{反物質}ではなく）と呼ぶのは、一部には慣習によるものである。物質とは、我々自身がそこから作られている方の対の成員である。物理的な主張は、観測可能な領域全体を通じて一方の符号が支配的であるというだけのことである。

\section{バリオン数と第一条件}

バリオン数は、陽子または中性子が\(B=1\)を持つように規格化される。クォークの言葉では、
\begin{equation}
B=\frac{N_q-N_{\bar q}}{3}.
\label{eq:c08-baryon-number}
\end{equation}
因子\(1/3\)は、各々の通常のバリオンが三つの価クォークを含むため本質的である。他のハドロンもバリオン数を持ちうるが、不安定なものは最終的に崩壊する。正味のクォーク超過がより一般的な記帳の道具である。

\begin{classroomqa}
\classquestion{正味のクォーク数は3で割るべきか、掛けるべきか。}
\lecturerresponse{3で割らなければならない。陽子は三つのクォークを持ちバリオン数は1なので、\(3B=N_q-N_{\bar q}\)である。黒板での訂正は定義の一部であり、物理的議論の変更ではない。\lecturetimestamp{00:09:23--00:10:17}}
\end{classroomqa}

バリオン数が電荷と似ているのは、可能な保存量としての側面だけである。電荷はまた長距離クーロン場の源でもあるが、バリオン数はそうではない。陽子は電荷をもつために電場をつくるが、中性子はバリオン数をもつものの正味の電荷をもたない。したがって、保存と力の生成は論理的に別個の性質である。

\begin{classroomqa}
\audiencequestion{バリオン数と電荷の間の関連する違いは何か。}
\lecturerresponse{どちらも帳簿上の電荷として機能しうるが、長距離電磁場の源となることが知られているのは電荷だけである。バリオン数は対応するクーロン力を生じない。\lecturetimestamp{00:11:03--00:12:15}}
\end{classroomqa}

もし\(B\)が厳密に保存されるなら、初期の過剰は永遠に凍結されたまま残るであろう。対生成と対消滅は\(N_B\)と\(N_{\bar B}\)を別々に変えうるが、等しい量だけしか変えられないので、それらの差は変わらないままである。したがって、釣り合った状態から出発する動的な理論は、バリオン数が変化しうるのでなければ成り立たない。

最も簡潔な模式的例は陽子崩壊である。電荷とエネルギーにより、我々は次を想像できる。
\begin{equation}
p\longrightarrow e^+ +\gamma,
\qquad
\bar p\longrightarrow e^-+\gamma.
\label{eq:c08-decay-pair}
\end{equation}
これらのチャネルはバリオン数破れの帳簿上の模型であり、支配的な実験チャネルの予言ではない。それらは電荷を保ちながら\(B=+1\)または\(-1\)をゼロにする。

\begin{figure}[htbp]
\centering
\begin{coursefitfigure}
\begin{tikzpicture}[x=1cm,y=1cm,>=Latex]
\draw[->,thick] (0,1.7)--(2.0,1.7);
\draw[->,thick] (2.0,1.7)--(4.1,2.55);
\draw[->,thick] (2.0,1.7)--(4.1,0.85);
\node[above] at (0.9,1.7) {\(p,\ B=1\)};
\node[right] at (4.1,2.55) {\(\gamma\)};
\node[right] at (4.1,0.85) {\(e^+\)};
\draw[->,thick] (6.0,1.7)--(8.0,1.7);
\draw[->,thick] (8.0,1.7)--(10.1,2.55);
\draw[->,thick] (8.0,1.7)--(10.1,0.85);
\node[above] at (6.9,1.7) {\(\bar p,\ B=-1\)};
\node[right] at (10.1,2.55) {\(\gamma\)};
\node[right] at (10.1,0.85) {\(e^-\)};
\node at (2.0,0.25) {\(B:1\to0\)};
\node at (8.0,0.25) {\(B:-1\to0\)};
\end{tikzpicture}
\end{coursefitfigure}
\caption{電荷共役な模式的チャネル。どちらの過程もバリオン数を破るが、等しい割合ならば優先される符号は生じない。}
\label{fig:c08-decay-redraw}
\end{figure}

\begin{classroomqa}
\audiencequestion{陽子は代わりに陽電子とニュートリノへ崩壊しうるか。}
\lecturerresponse{その二フェルミオン終状態は、一フェルミオン初期状態のフェルミオン parity に合わない。模式的な\(e^++\gamma\)終状態は一つのフェルミオンを含み、対称性の議論には十分である。現実的な陽子崩壊理論は他の粒子やチャネルを含みうるが、ここでの目的は\(\Delta B\ne0\)を示すことだけである。\lecturetimestamp{00:14:39--00:15:07}}
\end{classroomqa}

\begin{classroomqa}
\audiencequestion{厳密なバリオン数保存は対称性に対応するか。}
\lecturerresponse{はい。保存量は連続位相対称性と結びつけられうるが、その対称性に名前を与えるだけでは、自然がそれを尊重することを独立に証明するものではない。操作的検証はバリオン数を変える過程が起こるかどうかにとどまるので、対称性を保存則だけから推論するなら議論は循環的になる。\lecturetimestamp{00:15:08--00:15:46}}
\end{classroomqa}

\begin{classroomqa}
\audiencequestion{陽子--反陽子消滅それ自体を陽子崩壊と見なせるか。}
\lecturerresponse{用語は二次的であり、保存則はそうではない。\(p+\bar p\to\text{photons}\)では、初期バリオン数は\(1-1=0\)であり、終バリオン数もまたゼロである。消滅は正味の\(B\)を変えずに二つの粒子を取り除くが、式~\eqref{eq:c08-decay-pair}の一粒子チャネルは\(B\)を変える。\lecturetimestamp{00:15:48--00:16:47}}
\end{classroomqa}

現在の陽子は極めて安定である。講義は\(10^{32}\)--\(10^{34}\)年程度の上限を引用しており、宇宙の年齢よりはるかに長い。その観測は厳密な保存を確立するものではなく、現在のエネルギーではバリオン数を変える過程が極めて遅いことだけを教えている。

\section{離散対称性とバイアスの必要性}

次の問いは、法則が物質と反物質を区別するかどうかである。三つの離散変換が答えを整理する。
\begin{align}
C&:\ \text{particle}\longleftrightarrow\text{antiparticle},\\
P&:\ \mathbf x\longmapsto-\mathbf x,\\
T&:\ t\longmapsto-t.
\end{align}
ここで\(P\)は空間の handedness を反転し、\(T\)は微視的な映画を逆再生する。巨視的な第二法則は時間反転不変に見える必要はないが、問いはその背後にある微視的動力学に関するものである。

\begin{figure}[htbp]
\centering
\includegraphics[width=0.82\textwidth]{lecture_08_figure_03.png}
\caption{黒板には、修正された定義\(B=(N_q-N_{\bar q})/3\)が\(C\)、すなわち電荷共役の導入へと持ち込まれている。charge という語の後に続く未完成の書き込みは、独立の証拠としては用いられない。}
\label{fig:c08-charge-board}
\lectureframe{00:20:23.280}
\end{figure}

相対論的局所量子場理論は、\(C\)、\(P\)、\(T\)がそれぞれ別個に対称性であることを要求しない。その標準的な仮定の下では、合成変換\(CPT\)を要求する。ある過程が可能なら、すべての粒子をその反粒子に置き換え、空間配置を反映し、運動を逆転させると、別の許された過程が得られる。

量子力学における時間反転は、単に\(t\)の符号を変えるよりも微妙である。時間反転不変なポテンシャル中のスピンなし波動関数に対して、
\begin{equation}
i\hbar\frac{\partial\psi}{\partial t}=H\psi,
\qquad
T:\ \psi(\mathbf x,t)\longmapsto\psi^*(\mathbf x,-t).
\end{equation}
\(t\to-t\)を変えると時間微分が反転し、複素共役もまた\(i\to-i\)とするので、方程式が回復する。スピンをもつ粒子では、反ユニタリーな時間反転演算子は追加の構造を伴うが、複素共役はこの初等的な場合にすでに見えている。

\begin{classroomqa}
\audiencequestion{波動関数に対して、時間反転は\(t\to-t\)だけでなく複素共役も必要としますか。}
\lecturerresponse{はい。シュレーディンガー方程式がその理由を示している。時間微分だけを反転させると一方の辺だけが矛盾して変化するが、複素共役は\(i\)の符号も反転させる。\lecturetimestamp{00:23:58--00:25:04}}
\end{classroomqa}

ここで最初のサハロフ条件を述べることができる。宇宙が粒子と反粒子の等しい個数密度で始まったならば、何らかの過程が\(B\)を変化させない限り、対生成と対消滅は正味のバリオン数をゼロに保つであろう：
\begin{equation}
\boxed{\text{baryon-number violation}.}
\end{equation}
それがなければ、対称な初期状態が非ゼロのバリオン過剰を獲得することは決してできない。

陽子の現在の安定性はこの条件と両立する。大統一理論の推論では、バリオン数の破れは質量\(M_X\)の非常に重い粒子によって媒介されうる。\(M_X\)よりはるかに低いエネルギーでは、有効振幅は逆べき乗によって模式的に抑制される：
\begin{equation}
\begin{aligned}
\mathcal A_{\Delta B\ne0}&\propto \frac{E^2}{M_X^2},\\
\Gamma_p&\sim \frac{m_p^5}{M_X^4}.
\end{aligned}
\label{eq:c08-heavy-suppression}
\end{equation}
崩壊率の推定は、結合定数とチャネル依存因子を除いて、通常の次元6のスケーリングを持つ。非常に大きな\(M_X\)は、冷たい陽子を事実上永続的にする一方で、十分に高エネルギーの初期プラズマにおいてバリオン数を変える反応を可能にしうる。

\begin{classroomqa}
\audiencequestion{余分な衝突エネルギーはすでに重い質量\(M_X\)に含まれているのですか。}
\lecturerresponse{いいえ。\(M_X\)は媒介粒子の静止質量スケールであり、\(E\)は衝突で利用可能なエネルギーである。\(E/M_X\)における低エネルギー展開は\(E\ll M_X\)のとき強く抑制されるが、熱的衝突スケールが重い閾値に近づくにつれて抑制は弱まる。黒板に発見的に書かれた文字通りの分母は、この有効場の理論の意味で理解されるべきである。\lecturetimestamp{00:32:32--00:34:02}}
\end{classroomqa}

初期状態が粒子と反粒子の等しい個数密度を含んでいたとしよう。バリオン数の破れだけでは依然として偏りはない。すべての崩壊\(p\to e^++\gamma\)には荷電共役なパートナー\(\bar p\to e^-+\gamma\)があり、高温プラズマは逆反応も駆動するであろう。等しい微視的確率はバリオン数を通る流れを生むが、平均的なドリフトは生まない。

\begin{classroomqa}
\audiencequestion{反陽子のダイアグラムにおけるバリオン数の変化は何ですか。}
\lecturerresponse{入射反陽子は\(B=-1\)を持ち、電子と光子は\(B=0\)を持つので、このチャネルはバリオン数を\(-1\)からゼロに変化させる。その共役は\(+1\)をゼロに変化させる。二つのチャネルが等しく起こることは正味の選好を与えない。\lecturetimestamp{00:37:24--00:37:39}}
\end{classroomqa}

ランダムな揺らぎは、観測された過剰を説明するにはあまりにも小さすぎる。\(N\)個の独立な事象が二つの結果に対して等しい確率を持つならば、
\begin{equation}
\Delta N\sim\sqrt N,
\qquad
\frac{\Delta N}{N}\sim N^{-1/2}.
\end{equation}
講義の粗い見積もり\(N\sim10^{80}\)を観測可能な宇宙のバリオン数に用いると、
\begin{equation}
\frac{\Delta N}{N}\sim10^{-40},
\end{equation}
\(10^{-8}\)より32桁小さい。コイントスの揺らぎもまた、観測されるすべての領域で同じ符号が支配する理由を与えない。

\begin{classroomqa}
\audiencequestion{初期によりずっと大きな個数密度、おそらく我々の観測可能な宇宙になった共動領域に\(10^{160}\)対が存在し、それが\(10^{80}\)だけ揺らぐことはありえますか。}
\lecturerresponse{平方根の計算はそのような作為的な個数を許すが、観測された空間的コヒーレンスを説明しない。無制限の統計的揺らぎは自然に物質優勢と反物質優勢の領域を生むであろうし、可視宇宙全体にわたって一つの符号を生むことはない。\lecturetimestamp{00:39:38--00:40:18}}
\end{classroomqa}

講義時代の観測的議論は二種類の証拠を用いる。ヘリウムを含む通常の原子核は宇宙線中に存在するが、議論されているデータでは反原子核は観測されなかった。反陽子だけでは、高エネルギー衝突が二次的反陽子を生成しうるため、それほど決定的ではない。さらに、物質と反物質の銀河の境界は目立つ対消滅放射を生むであろう。必要な境界シグナルの不在は、観測可能な領域における銀河と反銀河の等しい混合を強く否定する。これらは講義の歴史的文脈における観測的議論であり、我々の地平線を超えた領域についての主張ではない。

\begin{classroomqa}
\audiencequestion{ヘリウム原子核自体が宇宙線中で測定されているのですか。}
\lecturerresponse{はい。荷電ヘリウム原子核は直接測定される。詳細な宇宙線統計は別の専門分野であるが、通常の原子核が一般的であることと、同程度の反原子核集団が存在しないことの対比がこの議論の要点である。\lecturetimestamp{00:41:53--00:42:22}}
\end{classroomqa}

したがって我々には微視的なバイアスが必要である。すなわち\(C\)の破れ、そして現実的なバリオン生成の議論では\(CP\)の破れである。中性\(B\)中間子系は実験室における一例を与える。中性中間子とその反粒子は、共役な終状態へ等しくない部分崩壊幅で崩壊しうる。
\begin{equation}
\Gamma(B^0\to f)\ne\Gamma(\bar B^0\to\bar f).
\label{eq:c08-cp-partial}
\end{equation}
厳密なチャネルと数値的な非対称性は板書では確定されなかった。頑健な教訓は、弱崩壊における粒子--反粒子非対称性の実験的に確立された存在である。

それにもかかわらず\(CPT\)定理は、真空中の粒子とその反粒子に対して等しい質量と等しい全崩壊幅を要求する。
\begin{equation}
\Gamma_{\rm tot}(p)=\Gamma_{\rm tot}(\bar p).
\end{equation}
これは各分岐比が一致することを強制しない。一対の共役チャネルにおける差は、模式的な電子チャネルやミューオンチャネルのような他のチャネルにおける差によって補償されうるが、和は等しいままである。

\begin{classroomqa}
\audiencequestion{陽子と反陽子の全崩壊率が等しいならば、どのようにしてバイアスが生じうるのか。}
\lecturerresponse{全幅の等しさはチャネルごとの等しさを意味しない。さらに、バリオン生成はプラズマ中で起こる。熱的占有因子、逆反応、散乱チャネルが運動方程式に入る。決定的な主張は、媒質が単に\(CPT\)を回避するということではなく、詳細釣り合いが失われれば\(CP\)非対称な部分過程がソース項を生成しうるということである。\lecturetimestamp{00:49:39--00:51:48; 00:53:22--00:57:14}}
\end{classroomqa}

これが第二のサハロフ条件である。
\begin{equation}
\boxed{C\ \text{and}\ CP\ \text{violation}.}
\end{equation}
これはバリオン数を変える反応に方向性を与える。サハロフの論文の直前に\(CP\)の破れが発見されたことで、この条件は物理的に具体的なものとなった。

最小標準模型は注意深い区別を要する。それは通常の摂動的な陽子崩壊を許さないが、その電弱アノマリーは\(B+L\)の非摂動的な破れを許す。高温では、スファレロン遷移がその破れを効率的にしうる。これが板書の後半の自己修正の訂正された内容である。標準模型のソースのみが宇宙の非対称性を定量的に再現するかどうかは、別の、はるかに難しい問題である。

\section{平衡と第三の条件}

なお一つの障害が残る。完全な熱平衡では、すべての反応はその逆反応によって釣り合っている。バリオン数が変化しうるならば、その平衡化学ポテンシャルは任意の初期\(B\)を保存しえず、詳細釣り合いが生成された過剰を消去する。講義は同じ点を\(CPT\)を通じて表現する。静的で平衡な熱浴には巨視的な時間の矢がなく、したがって\(T\)因子は過程とその逆過程の間の欠けている区別を与ええない。

実際の宇宙は膨張し冷却している。したがって巨視的な状態は前向き時間と後ろ向き時間を区別する。これが重要になるのは、膨張が微視的緩和と競合して成功する場合のみである。
\begin{equation}
\Gamma_{\rm reaction}\lesssim H.
\label{eq:c08-freezeout-condition}
\end{equation}
もし\(\Gamma_{\rm reaction}\gg H\)ならば、膨張にもかかわらず衝突が急速に平衡を回復する。反応率がハッブル率を下回れば、プラズマは再調整できず、存在量は凍結される。

\begin{classroomqa}
\audiencequestion{なぜ膨張だけでは自動的に系を平衡から外すのに十分でないのか。}
\lecturerresponse{非常にゆっくり膨張する熱浴は、微視的反応が絶えず再調整するため、任意に平衡に近いままである。離脱が有効になるのは、膨張または冷却の時間\(H^{-1}\)が関連する反応時間\(\Gamma^{-1}\)に比べて短い場合のみである。そのとき巨視的状態の映画はどちらが前向きかを明らかにする。\lecturetimestamp{01:01:22--01:02:26}}
\end{classroomqa}

三つのサハロフ条件をここでまとめることができる。
\begin{enumerate}
\item バリオン数の破れ、それにより\(B\)が変化しうる。
\item \(C\)および\(CP\)の破れ、それにより変化が好ましい符号をもつ。
\item 熱平衡からの離脱、それにより詳細釣り合いがその好みを打ち消せない。
\end{enumerate}
これらは議論中の機構にとって必要であり、過剰への構造的に十分な道筋を与える。具体的な高エネルギー模型、反応ネットワーク、宇宙論的歴史なしには、その大きさを決定しない。

\begin{classroomqa}
\audiencequestion{非対称性の測定値は高エネルギー理論を制約するか。}
\lecturerresponse{原理的には、はい。実際には、結果は多くの未知の質量、結合、位相、反応率、初期宇宙論の特徴に依存する。したがって、詳細な模型にコミットすることなしに、単一の測定値を狭い許容領域に変えることは難しい。\lecturetimestamp{01:04:20--01:04:52}}
\end{classroomqa}

\begin{classroomqa}
\audiencequestion{バリオン非対称性と電子非対称性はどのように結びついているのか。}
\lecturerresponse{電気的中性が両者を結びつける。この模式的なチャネルでは、陽子を取り除くと陽電子が生じ、バリオン不均衡を逆転させるとレプトン不均衡が補償する向きにずれる。より一般には、バリオン生成とレプトン生成は電荷保存と標準模型の反応を通じて結合しているが、保存される組み合わせは機構に依存する。\lecturetimestamp{01:04:54--01:05:59}}
\end{classroomqa}

\begin{classroomqa}
\audiencequestion{観測されている標準模型の\(CP\)の破れは、バリオン過剰を説明するには小さすぎるのか。}
\lecturerresponse{それが詳細な電気弱バリオン生成解析の標準的な結論であるが、講義は計算の模型依存性と高エネルギーにおける多くの未知を強調する。ここで擁護できる主張は、既知の低エネルギー\(CP\)の破れは存在するが、数値的に成功するバリオン生成模型には一般に追加の要素が必要である、というものである。\lecturetimestamp{01:06:01--01:07:50}}
\end{classroomqa}

\begin{classroomqa}
\audiencequestion{今日の電子は単に反陽子の消滅から生じたのか。}
\lecturerresponse{模式的なチャネルはそのような帳尻合わせを誘うが、初期のプラズマには独立でほぼ対称なレプトン集団も存在していた。関連する機構をバリオン生成またはレプトン生成として定式化し、その後の反応がセクター間で非対称性を移すこともできる。歴史的な語\emph{バリオン生成}を、すべての電子の一意な微視的系譜と誤解してはならない。\lecturetimestamp{01:07:51--01:08:44}}
\end{classroomqa}

\begin{classroomqa}
\audiencequestion{なぜ議論はすべてのレプトンではなく通常バリオンで枠組みされるのか。}
\lecturerresponse{バリオン物質は原子スペクトルやその他の天文観測を通じて目録化されるが、宇宙ニュートリノセクターを直接数えるのははるかに難しい。この命名は、観測的にアクセス可能であったものと、この主題の歴史を反映している。\lecturetimestamp{01:08:45--01:09:12}}
\end{classroomqa}

\section{消滅は過剰を残す}

高温では、対生成と対消滅は両方向に進む。
\begin{equation}
\gamma+\gamma\longleftrightarrow p+\bar p,
\qquad
E_{\gamma\gamma}\gtrsim2m_pc^2\approx2\,\mathrm{GeV}.
\end{equation}
浴が冷えるにつれ、バリオン対を生成するのに十分なエネルギーを持つ光子は指数関数的に希少になる。消滅はエネルギー的に許されたままだが、逆生成はもはや平衡を維持できない。ほとんどすべての陽子は反陽子と出会い、式~\eqref{eq:c08-baryon-asymmetry}の小さな過剰が残る。高エネルギーの消滅光子はすぐにそのエネルギーを熱浴と分かち合う。

\begin{classroomqa}
\audiencequestion{比\(10^{-8}\)は現在の数値なのか、また大規模な消滅はいつ起こったのか。}
\lecturerresponse{それはこの講義における現在のバリオン対光子の桁である。それ以前には、バリオンも反バリオンも多数存在していた。逆対生成が無効になると、それらの対称成分は消滅し、正味の過剰を残した。\lecturetimestamp{01:09:13--01:11:44}}
\end{classroomqa}

\begin{classroomqa}
\audiencequestion{陽子--反陽子消滅はいつ止まったのか。}
\lecturerresponse{反陽子が無視できる存在量まで枯渇したときに止まった。講義はこの出来事を大まかにしか位置づけない。GeVスケールの温度、すなわち元素合成より十分前、そして確実に再結合よりずっと前である。この議論から正確な時刻は推論されない。\lecturetimestamp{01:11:45--01:12:13}}
\end{classroomqa}

消滅光子は消えていない。それらは最終的に宇宙マイクロ波背景放射を供給する熱放射の歴史に属する。それでも放射が今日のエネルギー収支に寄与するのはわずかである。なぜなら
\begin{equation}
\rho_{\rm rad}\propto a^{-4},
\qquad
\rho_m\propto a^{-3}.
\end{equation}
光子数は多いが、各光子のエネルギーは赤方偏移し、物質はその静止質量エネルギーを保持するからである。

\begin{classroomqa}
\audiencequestion{消滅がそれほど多くの光子を作ったのなら、なぜ放射は今日力学的に無視できるのか。}
\lecturerresponse{光子は依然として存在し、主に宇宙マイクロ波背景放射として存在するが、そのエネルギー密度はスケール因子の余分な冪を失う。膨張はその数を希釈するとともに各光子を赤方偏移させる。\lecturetimestamp{01:12:14--01:13:18}}
\end{classroomqa}

\begin{classroomqa}
\audiencequestion{非常に重い媒介粒子は標準模型が不完全であることを意味するのか。}
\lecturerresponse{大統一陽子崩壊の媒介粒子は確かに最小標準模型を超える新物理であろう。標準模型の異常バリオン数破れに関する後の訂正は、そのような通常の陽子崩壊粒子を与えるものではなく、それ自体で完全な定量的バリオン生成模型を与えるものでもない。\lecturetimestamp{01:13:20--01:13:31}}
\end{classroomqa}

\begin{classroomqa}
\audiencequestion{二つの模式的な崩壊チャネルの間に大きな不等性があれば、初期に純粋な\(p\)-\(\bar p\)集団を、観測されている等しい電子数と陽子数に変換できるだろうか。}
\lecturerresponse{そのような帳尻合わせを考案することはできるが、熱平衡プラズマにはすでに到達可能なすべての粒子と反粒子の大きな集団が含まれている。さらに、観測されている\(CP\)非対称性は、一方のチャネルを恣意的に支配的にしてよいという許可証ではない。整合的な模型は、黒板の上の二本の矢印だけでなく、結合した反応ネットワーク全体を発展させなければならない。\lecturetimestamp{01:13:33--01:15:10}}
\end{classroomqa}

バリオン生成の議論が正直に終わるのはここである。定性的な材料は知られているが、観測されている数は導出されていない。より完全な答えは、より完全な粒子理論と、ごく初期宇宙の制御された歴史の両方を待たなければならない。

\section{第二の数：原始的なむら}

インフレーションは、宇宙マイクロ波背景放射によって定量的になったもう一つの問いに答える。なぜ宇宙は大きなスケールでこれほど一様で等方なのか。CMBは\(2.7\,\mathrm K\)付近の温度をもつ極めて精密な黒体スペクトルを示し、その温度はほぼすべての方向で同じである。かつては漠然とした宇宙論的な印象にすぎなかったものが、高精度の事実となった。

重力はこの滑らかさをいっそう驚くべきものにする。拡散はインクの一滴を広げ、その濃度を一様にする。引力としての重力は逆のことをする。過密な領域は周囲の物質を引き寄せ、その密度を増しながら隣接領域を空にする。したがって小さな摂動は、時間とともに成長することによって銀河、銀河団、ボイドの種となる。

\begin{figure}[htbp]
\centering
\begin{coursefitfigure}
\begin{tikzpicture}[x=1cm,y=1cm,>=Latex]
\draw[thick] (-4.2,0) ellipse (1.25 and 0.72);
\draw[very thick] (0,0) ellipse (1.55 and 0.92);
\draw[thick] (4.2,0) ellipse (1.25 and 0.72);
\node at (-4.2,0) {underdense};
\node at (0,0) {overdense};
\node at (4.2,0) {underdense};
\draw[->,thick] (-2.65,0)--(-1.65,0);
\draw[->,thick] (2.65,0)--(1.65,0);
\draw[->,thick] (-2.85,0.75)--(-1.45,0.4);
\draw[->,thick] (2.85,0.75)--(1.45,0.4);
\draw[->,thick] (-2.85,-0.75)--(-1.45,-0.4);
\draw[->,thick] (2.85,-0.75)--(1.45,-0.4);
\end{tikzpicture}
\end{coursefitfigure}
\caption{重力不安定性は、過小密度の周囲から過密領域へ物質を移し、初期のコントラストを拡散で消すのではなく増幅する。}
\label{fig:c08-gravity-growth}
\end{figure}

この成長を逆向きにたどると、はるかに滑らかな初期宇宙が示唆される。完全な一様性もまた受け入れられない。完全な対称性は持続し、銀河は形成されないからである。有用な尺度は密度コントラスト
\begin{equation}
\delta\equiv\frac{\delta\rho}{\rho},
\end{equation}
であり、ここで\(\delta\rho\)は背景密度からの局所的なずれである。ジム・ピーブルズや他の初期の宇宙論者に関連する構造成長の推定では、デカップリング時に原始的な振幅がおよそ
\begin{equation}
\boxed{\frac{\delta\rho}{\rho}\sim10^{-5}}
\end{equation}
であることが必要とされ、\(10^{-4}\)--\(10^{-5}\)が歴史的な範囲として挙げられていた。ここでもまた、その数は問いを鋭くするが、その値そのものをまだ説明してはいない。

1980年にアラン・グースによって導入された広義のインフレーション提案は、巨大でほぼ指数関数的な膨張の期間である。小さな滑らかな領域を何桁にもわたって引き伸ばすことで、既存の曲率を薄め、不規則性を観測可能な領域より大きなスケールへ押しやる。しわのある風船の類推は平滑化の直観をとらえているが、空間の文字どおりの埋め込みと受け取るべきではない。

\begin{classroomqa}
\audiencequestion{近くの物質が方向によって目に見えてむらだらけであるのに、どうして宇宙は等方でありうるのか。}
\lecturerresponse{等方性は大きなスケールでの統計的な主張であり、この部屋や銀河やあらゆる視線が同一に見えるという主張ではない。地球は局所的には山がちだが、その半径と比べれば極めて滑らかである。同様に、銀河や銀河団は、非常に滑らかな大規模背景の上に重なった局所構造である。\lecturetimestamp{01:25:37--01:26:22}}
\end{classroomqa}

\begin{classroomqa}
\audiencequestion{ランダムな過程なら、時には極端に過密な方向が生じるのではないか。}
\lecturerresponse{摂動は無制限の独立な揺らぎではない。それらは特定の、ほぼ一様で等方な統計的スペクトルをもつランダム場をなす。その振幅は、波長が大きくても\(10^{-5}\)程度と小さい。スペクトルとその起源は、恣意的に大きなコイントスの外れ値を仮定するものではなく、これから展開されるべき主題である。\lecturetimestamp{01:26:26--01:28:23}}
\end{classroomqa}

\begin{classroomqa}
\audiencequestion{インフレーションは初期時代に伴う高温によって引き起こされたのか。}
\lecturerresponse{インフレーションの間、温度は主要な変数ではない。中心的な現象は、場の応力エネルギーによって駆動される急速な指数関数的膨張である。熱浴が再び重要になるのは、インフレーションが終わり、そのエネルギーが粒子に変換されるときである。\lecturetimestamp{01:28:25--01:29:04}}
\end{classroomqa}

\section{力学的な前奏：抵抗を伴う運動}

インフレーションの方程式に備えて、単位質量の石が粘性流体中を落下する場合を考える。その鉛直座標を\(\phi\)と呼ぶ。まもなく石に取って代わるスカラー場を先取りしている。その位置エネルギーを\(V(\phi)\)とすれば、保存力は
\begin{equation}
F(\phi)=-\frac{dV}{d\phi},
\label{eq:c08-force-potential}
\end{equation}
であり、抵抗がなければニュートンの方程式は\(\ddot\phi=F\)である。

\begin{figure}[htbp]
\centering
\includegraphics[width=0.82\textwidth]{lecture_08_figure_06.png}
\caption{落下する石のアナロジーに対する、読み取り可能なボード上の関係\(F=-dV/d\phi\)。下向きの矢印は、\(\phi\)の増加が上向きを指すという規約を定めている。}
\label{fig:c08-force-board}
\lectureframe{01:35:09.513}
\end{figure}

線形粘性抵抗に対して、付加的な力は\(-\gamma\dot\phi\)であり、\(\gamma>0\)である。方程式は
\begin{equation}
\ddot\phi+\gamma\dot\phi=F(\phi)
=-V'(\phi).
\label{eq:c08-viscous-motion}
\end{equation}
となる。抵抗は静止時には消え、常に速度に逆らう。\(\phi\)が上向きに増加するなら、落下する石は\(F<0\)および\(\dot\phi<0\)を持ち、そのとき\(-\gamma\dot\phi>0\)であるから、抵抗は要求どおり上向きを指す。

\begin{classroomqa}
\classquestion{なぜポテンシャルによる力と抵抗の力は同じマイナス符号で現れるのか。それらは反対向きを指すべきではないのか。}
\lecturerresponse{それらの向きは異なる量によって決まる。ポテンシャルによる力は\(-V'(\phi)\)である。抵抗は\(-\gamma\dot\phi\)である。下向きの運動の間、\(F\)と\(\dot\phi\)はともに負であるから、\(F\)は下向きを指し、一方\(-\gamma\dot\phi\)は上向きを指す。したがって、代数はどちらの式も変えることなく、反対の物理的向きを生み出す。\lecturetimestamp{01:34:16--01:35:27}}
\end{classroomqa}

石が加速するにつれて、抵抗は力がつり合うまで増大する。\(F\)がほぼ一定であるとき、終端運動は
\begin{equation}
0=F-\gamma\dot\phi_{\rm term},
\qquad
\dot\phi_{\rm term}=\frac{F}{\gamma}.
\end{equation}
を満たす。符号は正しい。下向きの\(F\)は負の終端速度を与える。したがって、浅いポテンシャルと大きな\(\gamma\)は、坂を下る遅い運動を生み出す。

\begin{figure}[htbp]
\centering
\begin{tikzpicture}[x=1cm,y=1cm,>=Latex]
\draw[->] (0,0)--(7.0,0) node[right] {\(\phi\)};
\draw[->] (0,0)--(0,4.4) node[above] {\(V(\phi)\)};
\draw[very thick] plot[smooth] coordinates {(0.5,3.7) (1.6,3.55) (2.8,3.2) (4.0,2.55) (5.1,1.65) (6.3,0.6)};
\filldraw (1.45,3.58) circle (0.09);
\draw[->,thick] (1.55,3.48)--(2.25,3.23);
\node[align=left,anchor=west] at (2.45,3.75) {large drag,\\slow downhill motion};
\end{tikzpicture}
\caption{浅いポテンシャルと強い抵抗。場のアナロジーでは、石の座標は空間的に一様なスカラー振幅に置き換えられる。}
\label{fig:c08-shallow-potential}
\end{figure}

\section{インフラトンとハッブル摩擦}

インフレーションは、慣例的にインフラトンと呼ばれる付加的なスカラー場\(\phi\)を仮定する。歴史的には、これは宇宙の滑らかさを説明するために導入された大胆な仮説であった。まず、空間的にほぼ一様な場\(\phi=\phi(t)\)から始める。膨張は空間的変動を引き伸ばすが、一様な問題は本質的な力学を明らかにする。

正準スカラー場に対して、完全な局所エネルギー密度は
\begin{equation}
\rho_\phi
=
\frac12\dot\phi^{\,2}
+\frac{1}{2a^2}(\boldsymbol\nabla\phi)^2
+V(\phi).
\end{equation}
である。一様近似\(\boldsymbol\nabla\phi=0\)の下で、これは
\begin{equation}
\rho_\phi=\frac12\dot\phi^{\,2}+V(\phi).
\label{eq:c08-scalar-density}
\end{equation}
に帰着する。第一項は場の配位空間における運動エネルギーである。それは通常の空間を通る並進運動ではなく、時間的変化を測る。

\begin{classroomqa}
\audiencequestion{\(\phi\)のドットは時間に関する偏微分を意味するのか。}
\lecturerresponse{はい。一様な場に対して、\(\dot\phi=d\phi/dt\)である。空間勾配は別の正の項に寄与するが、現在の仮定の下ではそれらは消える。\lecturetimestamp{01:41:08--01:41:34}}
\end{classroomqa}

単位共動座標の箱を追跡する。その物理的体積は\(a^3(t)\)であるから、その内部のスカラーエネルギーは
\begin{equation}
E_\phi
=a^3(t)
\left[
\frac12\dot\phi^{\,2}+V(\phi)
\right].
\end{equation}
である。明示的な因子\(a^3(t)\)が新しい要素である。共動座標が固定されていても、物理的な箱は膨張する。

\begin{classroomqa}
\audiencequestion{スケール因子\(a\)は空間的位置に依存するのか。}
\lecturerresponse{一様なFRW模型では依存しない。与えられた宇宙時刻において空間全体で同じであるが、時間とともに変化しうる。もし\(\phi\)が位置に依存するなら、その空間勾配エネルギーを復活させなければならない。\lecturetimestamp{01:44:12--01:44:49}}
\end{classroomqa}

\begin{classroomqa}
\audiencequestion{今はただ一つの\(\phi(t)\)しか書かないが、場の値は空間の各点に関連付けられているのか。}
\lecturerresponse{はい。場は各時空点に値を割り当てる。一様性とは、それらの空間的な値が各時刻で一致することを意味し、力学を一つの共通関数\(\phi(t)\)に帰着させる。\lecturetimestamp{01:44:50--01:45:06}}
\end{classroomqa}

\begin{classroomqa}
\audiencequestion{場の運動項はどのような物理的運動を記述するのか。場は強くなっているのか、弱くなっているのか。}
\lecturerresponse{それは運動方程式、初期値と速度、そして\(V'(\phi)\)の符号によって決まる。アナロジーは、その運動の向きがエネルギー式だけからは分からない粒子である。\lecturetimestamp{01:45:07--01:46:08}}
\end{classroomqa}

一様なラグランジアンは運動エネルギーからポテンシャルエネルギーを引いたものであり、
\begin{equation}
L_\phi=a^3(t)
\left[
\frac12\dot\phi^{\,2}-V(\phi)
\right].
\label{eq:c08-scalar-lagrangian}
\end{equation}
オイラー--ラグランジュ方程式を適用すると
\begin{align}
\frac{d}{dt}\left(\frac{\partial L_\phi}{\partial\dot\phi}\right)
&=
\frac{\partial L_\phi}{\partial\phi},\\
\frac{d}{dt}\left(a^3\dot\phi\right)
&=-a^3V'(\phi).
\end{align}
が得られる。積の法則により二つの項が生じ、
\begin{equation}
a^3\ddot\phi+3a^2\dot a\,\dot\phi
=-a^3V'(\phi).
\end{equation}
\(a^3\)で割り、\(H=\dot a/a\)を定義すると、
\begin{equation}
\boxed{
\ddot\phi+3H\dot\phi+V'(\phi)=0
}
\label{eq:c08-hubble-friction}
\end{equation}
あるいは、\(F=-V'(\phi)\)として、
\begin{equation}
\ddot\phi+3H\dot\phi=F.
\end{equation}
が得られる。すべての行は式~\eqref{eq:c08-scalar-lagrangian}の膨張する体積因子から従う。

\begin{classroomqa}
\audiencequestion{\(H=\dot a/a\)はどこでも同じであり、本当に定数なのか。}
\lecturerresponse{一様な模型では、宇宙時刻のスライス上のすべての空間点で同じである。時間的に一定である必要はない。``ハッブル定数''はしばしばその現在値を指し、一方\(H(t)\)は進化を通じてのハッブル率である。\lecturetimestamp{01:50:32--01:50:57}}
\end{classroomqa}

式~\eqref{eq:c08-hubble-friction}はこの類推を精密なものにする。それは粘性運動と同じ形をしており、抗力係数は
\begin{equation}
\gamma_{\rm cosmic}=3H.
\end{equation}
である。したがって膨張は一様なスカラー場の時間発展を遅くする。\(H\)が大きく、\(V(\phi)\)が十分に平坦であれば、慣性項は抗力項に比べて小さくなり、場はスローロール関係
\begin{equation}
3H\dot\phi\simeq-V'(\phi).
\end{equation}
に近づく。冷たい粘性油の中を落下する球のように、場は浅いポテンシャルを横切るのに、たとえ力が働いていても、途方もなく長い時間を要しうる。

次の段階は、スカラーエネルギー密度をフリードマン方程式に代入し、膨張率をゆっくりと変化する場に応答させることである。この連成系がインフレーションを生み出す。今のところ、本質的な準備は整った。場は運動エネルギーとポテンシャルエネルギーを蓄え、膨張する体積はハッブル摩擦を生み出し、ゆっくりと転がる場はそのポテンシャルエネルギーをほぼ一定に保ち、急速な膨張を駆動するのに十分な長さ持続しうる。

\section{引き継ぐべきこと}

バリオン非対称性と原始密度ゆらぎは相補的な役割を果たす。\(10^{-8}\)も\(10^{-5}\)も、名前を付けるだけでは説明されないが、それぞれが漠然とした宇宙論的印象を定量的な問題へと変える。前者はサハロフの条件へと導き、後者はインフレーションへと導く。

インフレーションへの橋渡しも同様に精密である。FRW宇宙における一様なスカラー場は
\[
\ddot\phi+3H\dot\phi+V'(\phi)=0.
\]
に従う。中央の項は、挿入された現象論的粘性ではない。それは物理的体積\(a^3(t)\)の増大から導かれる。この一つの結果が、ポテンシャルに沿ってゆっくりと転がる場によって駆動される宇宙についての次回の講義の説明を組織するであろう。